\documentclass[11pt]{article}
\usepackage[margin=25mm]{geometry}
\usepackage{fontspec}
\setmainfont{Latin Modern Roman}
\usepackage{amsmath,amssymb,mathtools}
\usepackage{unicode-math}
\setmathfont{Latin Modern Math}
\usepackage{xurl,hyperref}
\hypersetup{colorlinks=true,urlcolor=blue}
\setcounter{secnumdepth}{0}
\setlength{\parindent}{0pt}
\setlength{\parskip}{5pt}
\setlength{\emergencystretch}{3em}
\newcommand{\sref}[2]{\hyperlink{source:#1:#2}{[S#2]}}
\newcommand{\cataloguescope}[1]{\paragraph{Recorded target.}#1}
\newcommand{\R}{\mathbb R}
\newcommand{\C}{\mathbb C}
\newcommand{\N}{\mathbb N}
\newcommand{\Z}{\mathbb Z}
\newcommand{\Q}{\mathbb Q}
\newcommand{\F}{\mathbb F}
\newcommand{\eps}{\varepsilon}
\begin{document}
% UnsolvedMath ID 20002530; source catalogue code AIM-PHYSICS-0038
\setcounter{section}{20002529}
\section{The signed Hill threshold for binary exchange}\label{20002530}
{\small UnsolvedMath ID 20002530 · Mathematical Physics}
\subsection{Definitions and mathematical statement}

\paragraph{Shared notation.}$\mathbb N=\{1,2,\ldots\}$; the standard number fields have their usual meanings. All additional conventions are stated locally.


\paragraph{Definitions and notation.}In the Newtonian $n$-body problem, positions $q_i\in\mathbb R^d$ and masses $m_i>0$ obey $m_i\ddot q_i=\nabla_{q_i}U$, where $U(q)=\sum_{i<j}m_im_j/|q_i-q_j|$ and the gravitational constant is normalized to one. Here $T=\tfrac12\sum_i m_i|\dot q_i|^2$, $E=T-U$, and a collision has $q_i=q_j$. The center of mass can be fixed at zero; moment of inertia is $I=\sum_i m_i|q_i|^2$. For a planar orbit write $C$ for scalar angular momentum; in spatial motion use its magnitude in $C^2E$. Let $L$ be the corresponding $C^2E$ value of the circular Euler collinear relative equilibrium for the stated masses. An exchange orbit connects one asymptotic binary-plus-escaping-body arrangement to another with a different binary. Hill stability here prohibits that exchange. A stable (respectively unstable) manifold consists of points tending to the invariant set in forward (respectively backward) time. A connecting chain must join the indicated invariant sets on compatible conserved-quantity levels.

\paragraph{Catalogue wording retained for traceability.}
For an exchange orbit in which $m_3$ escapes from the binary $\{m_1,m_2\}$ as $t\to-\infty$ and $m_2$ escapes from $\{m_1,m_3\}$ as $t\to+\infty$, how close to $1$ can $C^2E/L$ be, with $L=C^2E$ for circular Euler motion? The related question in the remarks asks whether a chain between the asymptotic binary invariant sets $P_{12},P_{13}$ can be constructed via the normally hyperbolic Euler invariant sphere $S^3_{\mathrm{Euler}}$, using suitable intersections of stable and unstable manifolds.
\paragraph{Scope and source qualification.}The quotient damaged by PDF extraction is repaired. The source records a negative-energy Hill-stability regime and suggests invariant-manifold connections; those are context, not a new proof or a universal exchange-existence assertion. The catalogue title/status was paired with a machine-generated research attempt; that attempt is not peer-reviewed evidence of settlement. The statement covers the original source questions rather than only the title’s subcase.
\subsection{Short English statement}
Find the sharp energy–angular-momentum threshold for three-body binary exchange.
\subsection{Research attempt}
\paragraph{Model, date and investigated scope.}
OpenAI GPT-6 Astra (configured model \texttt{gpt-6-astra}), ultra reasoning; generated 11 October 2026. Derived the amended-energy shape inequality and scaling/sign checks for C2E. No exchange orbit near the Euler threshold or invariant-manifold chain is constructed.

\paragraph{Sources and literature check.}
Read the supplied complete definition, catalogue wording, scope qualifications and source list for this canonical ID. The original-source citations are inherited as supplied references; no fresh primary-source or novelty search was performed for this attempt. No current open-status claim is inferred from this calculation.

\paragraph{An energy-angular-momentum necessary inequality.}Fix center of mass zero and let \(I=\sum m_i|q_i|^2\), \(T=\tfrac12\sum m_i|\dot q_i|^2\), \(U>0\). Cauchy--Schwarz for angular momentum gives \(C^2\le2IT\) in either the planar scalar or spatial-magnitude convention. Since \(E=T-U\),
\[E\ge C^2/(2I)-U.
\]
For a fixed shape with \(I=r^2\) and \(U=U_s/r\), \(U_s>0\), the right side is \(C^2/(2r^2)-U_s/r\). If \(C\ne0\), completing the square in \(1/r\) gives its minimum \(-U_s^2/(2C^2)\), hence every configuration on that shape obeys
\[C^2E\ge-U_s^2/2.
\]
Equality requires both the minimizing scale \(r=C^2/U_s\) and purely rotational velocity saturating the angular-momentum inequality. A rotating central configuration is the relevant dynamical situation, but the inequality alone does not prove it is reached.
\paragraph{Sign and scaling checks.}Newtonian dilation \(q(t)\mapsto a q(t/a^{3/2})\) sends \(E\mapsto E/a\), \(C\mapsto a^{1/2}C\), leaving \(C^2E\) invariant. When the Euler reference \(L<0\), dividing an inequality by \(L\) reverses its direction. The case \(C=0\) must be treated separately and cannot use the scale minimizer above.
\paragraph{Gap.}A necessary allowed-shape bound does not produce an exchange orbit or a chain through the Euler invariant sphere. The first missing step is proving the required stable/unstable intersections on matching conserved levels and following them without collision or escape into the wrong channel. No sharp approach to ratio one is established, and the sign-sensitive Hill-stability context is not promoted to a new sufficiency theorem.

\paragraph{Verification and adversarial audit.}
Checked Cauchy-Schwarz factors, scale minimum, equality requirements, C=0 exception, negative-L division and Newtonian dilation. This was an LLM adversarial audit, not independent human verification; no theorem prover was used. Prompts A then B applied to this ID: Checked Cauchy-Schwarz factors, scale minimum, equality requirements, C=0 exception, negative-L division and Newtonian dilation. Scope remains limited: Prove compatible stable/unstable manifold intersections and collision-free continuation through the desired exchange channels near the threshold.

Coordinator checked the angular-momentum bound, fixed-shape scale minimum and invariant Newtonian dilation product. This additional check was model-assisted, not independent human verification.

\paragraph{Tools, process and contributions.}
Python 3 and local shell used to read supplied files and serialize this manuscript; mathematics derived symbolically in the text, with no numerical experiment or formal proof assistant. Contribution prompt A was followed by prompt B. GPT-6 Astra with ultra reasoning drafted the research and performed its stated self-audit; any additional coordinator contribution is identified in the verification and audit. A GPT-6-based Codex coordinator assembled the manuscript and checked repository integration; its exact variant and reasoning configuration were not exposed. No human mathematical contribution, independent human proof check or formal proof certification is supplied. Complete derivations and the consulted source references are included here; private reasoning transcripts are not published.

\paragraph{Final status and first remaining gap.}
\textbf{UNRESOLVED}. Prove compatible stable/unstable manifold intersections and collision-free continuation through the desired exchange channels near the threshold.

\section{Completion Estimate}
\noindent\textbf{COMPLETION: 0\%}
The calculation is an elementary benchmark or conditional reduction; it establishes no claimed new portion of the complete source target.

\subsection{Sources}
\begin{itemize}
\item[\hypertarget{source:20002530:1}{S1}] ULAM.AI and the UnsolvedMath Contributors, UnsolvedMath, record 20002530 (AIM-PHYSICS-0038). Catalogue curation under CC BY 4.0; definitions and English exposition added locally.
\url{https://huggingface.co/datasets/ulamai/UnsolvedMath}.
\item[\hypertarget{source:20002530:2}{S2}] American Institute of Mathematics, Variational Methods in Celestial Mechanics; notes by Kuo-Chang Chen, version 22 June 2003, source item 1.
\url{https://aimath.org/WWN/varcelest/varcelest.pdf}.
\end{itemize}
\label{end:20002530}
\end{document}
